\title{DeepSeekMath: Pushing the Limits of Mathematical Reasoning in Open Language Models}

\begin{abstract}
Mathematical reasoning remains a substantial obstacle for language models owing to its intricacy and highly organized structure. In this work, we present DeepSeekMath~7B, developed by further pre-training DeepSeek-Coder-Base-v1.5 7B with 120B math-focused tokens derived from Common Crawl, along with supplementary natural language and programming code datasets. DeepSeekMath~7B attains a remarkable 51.7\% on the challenging, competition-oriented MATH benchmark, doing so independently of external tool integrations or ensembling strategies, and nearing the proficiency of models like Gemini-Ultra and GPT-4. With self-consistency evaluated over 64 completions, DeepSeekMath~7B reaches 60.9\% accuracy on the MATH task. The strength of DeepSeekMath~in mathematical reasoning arises from two principal elements: first, our design of a robust data curation framework that effectively leverages the breadth of publicly accessible web resources; second, the development of Group Relative Policy Optimization (GRPO), an adaptation of Proximal Policy Optimization (PPO), which not only enhances mathematical reasoning, but also improves PPO’s memory efficiency.
\end{abstract}

\section{Introduction}

Large language models (LLM) have transformed artificial intelligence methods for mathematical reasoning, leading to notable progress on quantitative reasoning benchmarks \citep{MATH} as well as geometry reasoning benchmarks \citep{trinh2024solving}. Additionally, these models have demonstrated considerable value in supporting humans tackling challenging mathematical tasks \citep{tao}. Nevertheless, state-of-the-art systems like GPT-4 \citep{gpt4} and Gemini-Ultra \citep{gemini} are not open to the public, while open-source alternatives presently exhibit much lower performance.

This work presents DeepSeekMath, a specialized language model tailored for mathematics, which substantially exceeds the performance of current open-source models and narrows the gap with GPT-4 on rigorous academic benchmarks. Central to this advance is the DeepSeekMath Corpus: an extensive, meticulously curated dataset of 120B math-related tokens used for pre-training. The collection process draws from Common Crawl (CC), applying a fastText-based classifier \citep{joulin2016fasttext}. For the classifier’s initial round, samples from OpenWebMath \citep{openwebmath} provide positive examples, while a varied array of unrelated web pages constitutes the negative set. The classifier then identifies further positive samples from the CC; these are subjected to detailed human annotation for higher quality, and the revised data is leveraged to further refine the classifier. Evaluation metrics confirm the corpus’s robustness: the DeepSeekMath-Base 7B foundation model secures a 64.2\% score on GSM8K \citep{gsm8k} and 36.2\% on the MATH dataset \citep{MATH}—results that surpass those achieved by Minerva 540B \citep{minerva}. Owing to the multilingual nature of the DeepSeekMath Corpus, we also observe gains on Chinese mathematical assessment tasks \citep{wei2023cmath,agieval}. We propose that our methodology in constructing and processing mathematical data lays the groundwork for further scholarly innovation, with ample prospects for future enhancement.

DeepSeekMath-Base is built upon DeepSeek-Coder-Base-v1.5 7B \citep{deepseek-coder}, as empirical evidence suggests that models pre-trained on code yield superior performance for this purpose over those with general LLM pre-training. Additionally, our experiments reveal that mathematical fine-tuning leads to improvements on benchmarks such as MMLU \citep{mmlu} and BBH \citep{bbh}, demonstrating not only gains in mathematical proficiency but also enhanced general reasoning skills.

Following the pre-training phase, we conduct mathematical instruction tuning on DeepSeekMath-Base utilizing datasets that include chain-of-thought \citep{cot}, program-of-thought \citep{pot,pal}, and tool-integrated reasoning \citep{tora}. The final model, DeepSeekMath-Instruct 7B, surpasses all other 7B models and demonstrates performance on par with leading 70B open-source models tuned for instructions.

In addition, we present Group Relative Policy Optimization (GRPO), an alternative reinforcement learning (RL) approach derived from Proximal Policy Optimization (PPO) \citep{schulman2017proximal}. Distinct from conventional methods, GRPO eliminates the need for a critic network and instead uses group scores to compute baselines, leading to significant reductions in computational requirements. Utilizing only a subset of English instruction tuning samples, GRPO achieves notable gains over the robust DeepSeekMath-Instruct across a range of tasks, demonstrating improvements on both in-domain benchmarks (GSM8K: 82.9\% $\rightarrow$ 88.2\%, MATH: 46.8\% $\rightarrow$ 51.7\%) and out-of-domain problems (e.g., CMATH: 84.6\% $\rightarrow$ 88.8\%) during RL training. Additionally, we propose a unified framework that encompasses existing methods like Rejection Sampling Fine-Tuning (RFT) \citep{yuan2023scaling}, Direct Preference Optimization (DPO) \citep{dpo}, PPO, and our GRPO. Through this unified framework, we observe that these approaches can be understood as either direct implementations of RL or as simplifications thereof. We further perform a series of systematic investigations—including analyses of online versus offline training, comparisons of outcome versus process supervision, and explorations of single-turn versus iterative RL—to better elucidate the critical components of this framework. Finally, we discuss how our RL methodology enhances the performance of instruction-tuned models and outline promising avenues for advancing RL effectiveness within this unified perspective.

\subsection{Contributions}
We present a scalable approach to math pre-training and undertake a comprehensive investigation and assessment of reinforcement learning methodologies.

\textbf{Math Pre-Training at Scale}

\begin{itemize}[topsep=0pt]
    \item Our study demonstrates that Common Crawl's openly available data holds significant utility for mathematical applications. Through the deployment of a carefully crafted data selection process, we assembled the DeepSeekMath Corpus—a high-quality collection comprising 120B tokens sourced from web pages curated for their mathematical relevance. This dataset is nearly seven times larger than the collection of math-related web pages used in Minerva \citep{minerva}, and approximately nine times greater in size than the newly introduced OpenWebMath \citep{openwebmath}.

\item The DeepSeekMath-Base 7B model, after pre-training, demonstrates performance on par with Minerva 540B \citep{minerva}. This suggests that parameter count is not the sole determinant of proficiency in mathematical reasoning; smaller models trained with high-quality datasets can also attain impressive results.

\item We present insights gathered from our experiments involving mathematical training. Engaging in code training before mathematical instruction enhances the models’ capacity to tackle mathematical tasks, regardless of whether tools are used. This provides evidence towards resolving the enduring inquiry: \textit{does code training enhance reasoning skills?} Our results indicate that it does, specifically in the domain of mathematical reasoning.

\item While it is customary to train models using arXiv papers, particularly for mathematics-focused tasks, this practice yields no significant gains across any of the mathematical benchmarks evaluated in the present study.
\end{itemize}

\textbf{Exploration and Analysis of Reinforcement Learning}

\begin{itemize}[topsep=0pt]
    \item We present Group Relative Policy Optimization (GRPO), a reinforcement learning algorithm characterized by its simplicity and efficiency. Unlike traditional methods such as Proximal Policy Optimization (PPO), GRPO eliminates the need for a critic model by utilizing group scores to estimate the baseline, thereby considerably reducing the computational demands during training.
    \item Our results show that the application of GRPO leads to notable improvements in the instruction-tuned model DeepSeekMath-Instruct, employing only instruction-tuning data. Additionally, we note that GRPO yields better generalization to out-of-domain tasks in the course of the reinforcement learning stage.
    \item We establish a comprehensive framework for understanding a range of algorithms, including RFT, DPO, PPO, and GRPO. In addition, we perform a thorough experimental analysis—investigating dimensions such as online versus offline training, outcome versus process supervision, and the effects of single-turn compared to iterative reinforcement learning—to dissect the fundamental aspects underpinning this framework.
    \item Leveraging our unified paradigm, we investigate the mechanisms contributing to the efficacy of reinforcement learning, and outline several prospective pathways toward further enhancing reinforcement learning methodologies for LLMs.
\end{itemize}

\subsection{Summary of Evaluations and Metrics}

\begin{itemize}[topsep=0pt]
    \item \textbf{English and Chinese Mathematical Reasoning}: We perform thorough evaluations of our models on both English and Chinese datasets, spanning mathematical questions from elementary to collegiate levels. The English datasets consist of GSM8K \citep{gsm8k}, MATH \citep{MATH}, SAT \citep{llemma}, OCW Courses \citep{minerva}, and MMLU-STEM \citep{mmlu}. For Chinese, the benchmarks include MGSM-zh \citep{mgsm}, CMATH \citep{wei2023cmath}, Gaokao-MathCloze \citep{agieval}, and Gaokao-MathQA \citep{agieval}. We assess two key capabilities of the models: generating complete text-based solutions without reliance on external tools, and solving questions through Python programming.

On English evaluation tasks, DeepSeekMath-Base matches the performance of the proprietary Minerva 540B model \citep{minerva} and outperforms all other open-source base models, such as Mistral 7B \citep{mistral} and Llemma-34B \citep{llemma}, irrespective of whether those models have benefited from math-specific pre-training, frequently achieving a substantial lead. Importantly, DeepSeekMath-Base demonstrates even stronger results on Chinese assessment datasets, which is plausibly attributable to our distinct pre-training approach that, unlike previous efforts \citep{minerva,llemma}, does not limit training data to English-language mathematical content but incorporates high-quality datasets in other languages as well. Through the application of mathematical instruction tuning and reinforcement learning, our advanced models, DeepSeekMath-Instruct and DeepSeekMath-RL, achieve a remarkable milestone by exceeding 50\% accuracy on the challenging, competition-level MATH benchmark—a first within the open-source community.

\item \textbf{Formal Mathematics}: We assess DeepSeekMath-Base on the informal-to-formal theorem proving benchmark introduced in \citep{dsp_proof}, utilizing the miniF2F dataset \citep{minif2f} with Isabelle \citep{isabelle} serving as the selected proof assistant. DeepSeekMath-Base achieves notable performance in few-shot autoformalization tasks.
\item \textbf{Natural Language Understanding, Reasoning, and Code}: In order to obtain a holistic understanding of the model's proficiency in general language comprehension, reasoning, and programming, we benchmark DeepSeekMath-Base against Massive Multitask Language Understanding (MMLU) \citep{mmlu}, featuring 57 multiple-choice problems spanning various fields, as well as BIG-Bench Hard (BBH) \citep{bbh}, comprised of 23 tasks predominantly requiring multi-step inference. Additionally, HumanEval \citep{codex} and MBPP \citep{mbpp} are used as standard benchmarks for code generation. The math-specific pre-training enhances model performance in both linguistic reasoning and understanding.

\section{Math Pre-Training}

\subsection{Data Collection and Decontamination} This section details the methodology employed to build the DeepSeekMath Corpus using data sourced from Common Crawl. As illustrated in Figure \ref{fig:data_collect}, an iterative workflow is introduced, which enables the structured extraction of a comprehensive mathematical corpus from Common Crawl, beginning with a seed dataset (such as a limited, yet high-quality, math-specific data collection). Importantly, this methodology can be generalized and adapted to additional fields, including but not limited to coding.

Initially, we select OpenWebMath \citep{openwebmath}, which comprises a curated set of mathematical texts sourced from the web, as our foundational seed corpus. We leverage this corpus to train a fastText model \citep{joulin2016fasttext}, which is subsequently used to identify additional web pages exhibiting characteristics similar to OpenWebMath. For training, we draw a random sample of 500,000 items from the seed corpus as positive examples, while a separate 500,000 web pages from Common Crawl serve as negative instances. An open-source fastText library\footnote{\url{https://fasttext.cc}} is utilized for training the model, with specific settings: a 256-dimensional embedding space, a learning rate set at 0.1, word n-grams up to length 3, a minimum word frequency threshold of 3, and 3 training epochs. To downsize the raw Common Crawl data, we implement both URL-based deduplication and near-deduplication strategies, reducing the corpus to 40B unique HTML web documents. Using the trained fastText model, we retrieve mathematical web content from this deduplicated set. To eliminate low-quality mathematical material, we employ fastText-predicted scores to rank these pages, retaining only those with the highest rankings. We evaluate the quantity of the filtered data by running pre-training on datasets comprising the top 40B, 80B, 120B, and 160B tokens. For this initial iteration, we opt to preserve only the top 40B tokens.

Following the initial round of data gathering, a significant number of mathematical web pages are still excluded, primarily due to the limited diversity in the positive training examples used by the fastText model. To address this, we expand the range of mathematical sources within the seed corpus, aiming to improve the fastText model’s performance. The process begins by partitioning the entire Common Crawl dataset according to domain, where a domain encompasses all pages sharing the same base URL. For each domain, we determine the proportion of pages retrieved during the first data collection cycle. Domains in which over 10\% of pages were captured are categorized as math-oriented (e.g., \url{mathoverflow.net}).

Next, we manually review and annotate URLs that specifically link to mathematical content within these math-related domains (e.g., \url{mathoverflow.net/questions}). Any relevant web pages corresponding to these URLs that were overlooked in prior iterations are subsequently incorporated into the seed corpus. This methodology enhances the collection of positive samples, thereby yielding a more effective fastText model that can extract a greater quantity of mathematical data in ensuing rounds. After completing four cycles of data collection, the resulting dataset consists of 35.5M mathematical web pages, amounting to 120B tokens in total. Upon analyzing results from the fourth round, we observe that nearly 98\% of the dataset had already been recovered in the third round, leading us to terminate further data collection at that point.

In order to prevent benchmark contamination, we adopt the filtering strategy proposed in \cite{deepseek-coder}, excluding any web content featuring questions or answers from English mathematical benchmarks, including GSM8K \citep{gsm8k} and MATH \citep{MATH}, as well as Chinese benchmarks like CMATH \citep{wei2023cmath} and AGIEval \citep{agieval}. Specifically, our filtering procedure eliminates any passage that contains a contiguous 10-gram sequence precisely matching any substring found in the benchmark datasets. For benchmark passages shorter than 10 grams but at least 3 grams long, we apply an exact substring match to identify and discard potentially contaminated data.

\subsection{Validating the Quality of the DeepSeekMath~Corpus}

We run pre-training experiments to investigate how the DeepSeekMath~Corpus is compared with the recently released math-training corpora:

\begin{itemize}[topsep=0pt]
    \item \textbf{MathPile} \citep{mathpile}: a dataset compiled from multiple sources, comprising 8.9B tokens, which includes materials from textbooks, Wikipedia, ProofWiki, CommonCrawl, StackExchange, and arXiv, with more than 85\% of the content originating from arXiv;
    \item \textbf{OpenWebMath} \citep{openwebmath}: a 13.6B-token subset of CommonCrawl, curated by filtering specifically for mathematical content;
    \item \textbf{Proof-Pile-2} \citep{llemma}: this mathematical dataset brings together OpenWebMath, AlgebraicStack (10.3B tokens representing math-focused code), and arXiv articles (28.0B tokens). In our experiments utilizing Proof-Pile-2, we adhere to the proportion of arXiv:Web:Code as 2:4:1, following the protocol in \cite{llemma}.
\end{itemize}

\subsubsection{Training Setting}
\label{sec:quality-policy}
For our experiments, we fine-tune a general-purpose language model with 1.3B parameters that shares its architecture with the DeepSeek LLMs \citep{deepseek-llm}; we refer to this model as DeepSeek-LLM 1.3B. Each mathematical dataset is used independently to train a separate model for 150B tokens. The entirety of the training process utilizes the HAI-LLM framework \citep{haillm}, which is designed for efficiency and minimal resource usage. Adhering to DeepSeek LLMs’ established procedures, we employ the AdamW optimizer \citep{adamW} with hyperparameters set to $\beta_1=0.9$, $\beta_2=0.95$, and $\mathrm{weight\_decay}=0.1$. The learning rate is managed using a multi-phase schedule: it linearly increases and reaches its highest point after an initial 2,000-step warmup, drops to 31.6\% of this maximum by 80\% of total training, and then is further reduced to 10.0\% by the 90\% mark. The learning rate caps at 5.3e-4. During training, a batch size comprising 4M tokens and a context window of 4K tokens are used.

\subsubsection{Evaluation Results}

\textbf{The DeepSeekMath~Corpus is of high quality, covers multilingual mathematical content, and is the largest in size.}

\begin{itemize}[topsep=0pt]
    \item \textbf{High-quality}:
    To assess downstream efficacy, we conduct few-shot chain-of-thought prompting \cite{cot} across 8 mathematical benchmarks. Table \ref{tab:corpora_comparison} reveals that the DeepSeekMath~Corpus consistently supports superior model performance compared to alternatives. Furthermore, Figure \ref{fig:corpora_comparisons} highlights that a model trained on DeepSeekMath achieves higher scores than one trained on Proof-Pile-2 with 50B tokens (equivalent to a single epoch over Proof-Pile-2), suggesting the overall quality of data in DeepSeekMath Corpus is distinctly higher.
    \item \textbf{Multilingual}:
    DeepSeekMath~Corpus contains diverse linguistic data, with English and Chinese comprising the majority. As indicated in Table \ref{tab:corpora_comparison}, utilizing DeepSeekMath~Corpus for training yields notable improvements in mathematical reasoning for both English and Chinese language tasks. By contrast, most prior mathematical corpora, focused chiefly on English, produce marginal gains or even decrease proficiency when applied to Chinese mathematical reasoning scenarios.
    \item \textbf{Large-scale}:
    In terms of scale, DeepSeekMath~Corpus is markedly larger than other established mathematical corpora. As shown in Figure \ref{fig:corpora_comparisons}, DeepSeek-LLM 1.3B, after being trained on DeepSeekMath~Corpus, experiences a sharper and more sustained advancement in learning performance. By comparison, conventional corpora are significantly smaller in size, necessitating multiple reuse cycles throughout training, which in turn leads to models quickly reaching a plateau in terms of performance gains.
\end{itemize}

\subsection{Training and Evaluating DeepSeekMath-Base 7B}

This section presents DeepSeekMath-Base 7B, a foundational model that demonstrates notable reasoning skills, particularly in mathematical contexts. The model is built upon DeepSeek-Coder-Base-v1.5 7B \citep{deepseek-coder} and undergoes training on 500B tokens. The dataset composition is as follows: 56\% originates from the DeepSeekMath Corpus, 4\% from AlgebraicStack, 10\% from arXiv, 20\% comprises code from Github, and the final 10\% consists of English and Chinese natural language data extracted from Common Crawl. The training protocol largely aligns with the setup detailed in Section \ref{sec:quality-policy}, with modifications including setting the peak learning rate to 4.2e-4 and employing a batch size of 10M tokens.

This study presents an extensive evaluation of DeepSeekMath-Base 7B’s mathematical proficiency, emphasizing its performance in independently generating complete mathematical solutions, employing external tools for mathematical problem-solving, and executing formal theorem-proving tasks. In addition to mathematics, we offer a broader characterization of the base model by examining its capabilities in natural language understanding, logical reasoning, and programming.

\paragraph{Mathematical Problem Solving with Step-by-Step Reasoning}

We assess the capabilities of DeepSeekMath-Base in addressing mathematical questions through the use of few-shot chain-of-thought prompting \citep{cot}, utilizing eight distinct benchmarks in both English and Chinese. These evaluation sets include benchmarks aimed at quantitative reasoning (such as GSM8K \citep{gsm8k}, MATH \citep{MATH}, and CMATH \citep{wei2023cmath}), as well as multiple-choice assessments (including MMLU-STEM \citep{mmlu} and Gaokao-MathQA \citep{agieval}). Together, they span a wide spectrum of mathematical domains, ranging from basic to advanced (elementary through college-level) topics.

Table \ref{tab:base_math_cot} illustrates that among open-source base models, DeepSeekMath-Base 7B demonstrates the strongest results across all eight evaluated benchmarks—including Mistral 7B \citep{mistral}, a commonly adopted general-purpose model, and Llemma 34B \citep{llemma}, a recent entrant trained on Proof-Pile-2 \citep{llemma} for mathematical tasks. In particular, DeepSeekMath-Base achieves more than a 10\% absolute improvement over all open-source base models on the MATH dataset, which assesses competition-level mathematical reasoning. Additionally, it outperforms Minerva 540B \citep{minerva}, a proprietary base model with 77 times the parameter count that is based on PaLM \citep{palm} and additionally pre-trained on math-specific data.

\paragraph{Mathematical Problem Solving with Tool Use} We assess the capability of models to perform mathematical reasoning assisted by external tools on the GSM8K and MATH benchmarks, employing few-shot program-of-thought prompting as outlined in \citep{pot,pal}. For each task, models are instructed to generate a Python script, taking advantage of libraries like \textit{math} and \textit{sympy} for advanced or complex calculations. The final answer is determined by running the generated code and using its output. According to Table \ref{tab:base_math_pot_proof}, DeepSeekMath-Base 7B achieves better performance than the previously best-performing model, Llemma 34B.

\paragraph{Formal Mathematics}

The automation of formal proofs plays an increasingly prominent role in ensuring the correctness and dependability of mathematical arguments, while also improving productivity. In this study, we assess DeepSeekMath-Base 7B on the informal-to-formal proof generation task as presented in \citep{dsp_proof}. This task involves generating a formal proof from an informal statement, its corresponding formalized version, and an informal proof. Our evaluation utilizes the miniF2F benchmark \citep{minif2f}, which focuses on Olympiad-level mathematical problems formalized in Isabelle. For each problem, we prompt the model with a small number of examples to produce a formal proof. Consistent with the methodology in \cite{dsp_proof}, our approach employs language models to generate proof outlines, subsequently relying on the standard automated prover Sledgehammer \citep{sledgehammer} to fill in unresolved steps. Table \ref{tab:base_math_pot_proof} illustrates that DeepSeekMath-Base 7B achieves notable results in the automatic formalization of mathematical proofs.

\paragraph{Natural Language Understanding, Reasoning, and Code}

We assess the natural language comprehension of models using MMLU \citep{mmlu}, their reasoning abilities through BBH \citep{bbh}, and their coding skills with the HumanEval \citep{codex} and MBPP \citep{mbpp} benchmarks. Table \ref{tab:understanding_reasoning_code} demonstrates that DeepSeekMath-Base 7B achieves marked improvements on MMLU and BBH in comparison to its predecessor, DeepSeek-Coder-Base-v1.5 \citep{deepseek-coder}, underscoring the beneficial effects of math-focused training on both understanding and reasoning capabilities. Furthermore, through the inclusion of code tokens during continual training, DeepSeekMath-Base 7B is able to preserve the coding performance level of DeepSeek-Coder-Base-v1.5 across both coding assessments. In summary, DeepSeekMath-Base 7B consistently surpasses the more general Mistral 7B model \citep{mistral} on all three reasoning and coding evaluation tasks.

\section{Supervised Fine-Tuning}

\subsection{SFT Data Curation}

We develop a mathematically focused instruction-tuning dataset that spans both English and Chinese problems, drawn from diverse branches of mathematics and featuring a range of difficulty levels. Each problem is accompanied by a corresponding solution in formats such as chain-of-thought (CoT) \citep{cot}, program-of-thought (PoT) \citep{pot,pal}, and reasoning with tool integration \citep{tora}. This dataset consists of a total of 776K training samples.

\begin{itemize}[topsep=0pt]
    \item \textbf{English mathematical datasets}: Our annotations extend to GSM8K and MATH problem sets, providing solutions that incorporate tool-based reasoning. Additionally, we utilize a curated subset of MathInstruct \citep{MathInstruct} and the Lila-OOD training set \citep{lila}, both featuring problems addressed through CoT or PoT methodologies. The English dataset encompasses a broad array of mathematical areas, including algebra, probability, number theory, calculus, and geometry.
    \item \textbf{Chinese mathematical datasets}: Our compilation consists of K-12 mathematical problems in Chinese, covering 76 distinct sub-topics (such as linear equations). Solutions are provided using both CoT and tool-integrated reasoning strategies.
\end{itemize}

\subsection{Training and Evaluating DeepSeekMath-Instruct 7B}

This section presents DeepSeekMath-Instruct 7B, which is derived from DeepSeekMath-Base by applying instruction tuning tailored for mathematical tasks. Training data are concatenated in random order until the cumulative sequence reaches a 4K token context window. The model is then trained for 500 steps, utilizing a batch size of 256 and maintaining a fixed learning rate of 5e-5 throughout the training process.

We evaluate models' mathematical performance both without and with tool use, on 4 quantitative reasoning benchmarks in English and Chinese. We benchmark our model against the leading models of the time:

\begin{itemize}[topsep=0pt]
    \item \textbf{Closed-source models} encompass: (1) the GPT series, with GPT-4 \citep{gpt4} and GPT-4 Code Interpreter~\footnote{\url{https://openai.com/blog/chatgpt-plugins\#\#code-interpreter}} representing the highest-performing variants, (2) Gemini Ultra and Gemini Pro \citep{gemini}, (3) Inflection-2 \citep{inflection-2}, (4) Grok-1~\footnote{\url{https://x.ai/model-card}}, along with recently launched models from Chinese firms such as (5) Baichuan-3~\footnote{\url{https://www.baichuan-ai.com}}, and (6) the most recent GLM-4~\footnote{\url{https://open.bigmodel.cn/dev/api\#glm-4}}

    from the GLM family \citep{glm}.

These models serve broad applications and have typically been refined through various alignment processes.
    \item \textbf{Open-source models} feature: general-purpose architectures such as (1) DeepSeek-LLM-Chat 67B \citep{deepseek-llm}, (2) Qwen 72B \citep{qwen}, (3) SeaLLM-v2 7B \citep{seallm}, and (4) ChatGLM3 6B \citep{chatglm3}. Additionally, specialized models with mathematical capability enhancements include: (5) InternLM2-Math 20B~\footnote{\url{https://github.com/InternLM/InternLM-Math}}, which extends InternLM2 by incorporating mathematical corpus training and subsequent instruction tuning; (6) Math-Shepherd-Mistral 7B, trained with PPO \citep{schulman2017proximal} leveraging a process-supervised reward model built on top of Mistral 7B \citep{mistral}; (7) the WizardMath collection \citep{wizardmath}, designed to enhance Mistral 7B and Llama-2 70B \citep{llama2} for mathematical reasoning using evolve-instruct methodology (which adapts instruction tuning with AI-generated instructions) and PPO optimization, predominantly utilizing GSM8K and MATH datasets; (8) MetaMath 70B \citep{metamath}, developed by further fine-tuning Llama-2 70B on expanded GSM8K and MATH data; (9) ToRA 34B \cite{tora}, which is a CodeLlama 34B variant tailored for tool-augmented mathematical reasoning; and (10) MAmmoTH 70B \citep{MathInstruct}, produced through instruction tuning of Llama-2 70B on MathInstruct.

Referencing Table \ref{tab:sft_rl_math}, when evaluated in scenarios without access to tool use, DeepSeekMath-Instruct 7B exhibits robust capabilities in step-by-step mathematical reasoning. On the competition-level MATH benchmark, our model exceeds the performance of every open-source alternative and outperforms most commercial solutions (such as Inflection-2 and Gemini Pro) by a minimum margin of 9\% absolute. This performance advantage holds even when compared to much larger systems (like Qwen 72B) or those refined via specialized math-oriented reinforcement learning approaches (e.g., WizardMath-v1.1 7B). Although DeepSeekMath-Instruct matches the performance of proprietary Chinese models such as GLM-4 and Baichuan-3 on this dataset, it still falls short of the leading results achieved by GPT-4 and Gemini Ultra.

When assessed in a scenario that permits both natural language-based reasoning and the utilization of programmatic tools for solving tasks, DeepSeekMath-Instruct 7B attains nearly 60\% accuracy on MATH, outperforming all prior open-source models. Across additional benchmarks, our model exhibits performance on par with DeepSeek-LLM-Chat 67B, the previous state-of-the-art, despite having only a tenth of its parameter count.
\section{Reinforcement Learning}

\subsection{Group Relative Policy Optimization} Reinforcement learning (RL) has demonstrated considerable success in enhancing the mathematical reasoning skills of LLMs following the Supervised Fine-Tuning (SFT) phase \citep{wang2023math,wizardmath}. In this part, we present Group Relative Policy Optimization (GRPO), our novel and resource-efficient RL method.

\subsubsection{From PPO to GRPO} Proximal Policy Optimization (PPO) \citep{schulman2017proximal} is an actor-critic RL algorithm that is widely used in the RL fine-tuning stage of LLMs \citep{ouyang2022training}. In particular, it optimizes LLMs by maximizing the following surrogate objective:

\begin{equation}
\footnotesize
    \mathcal{J}_{PPO}(\theta) = \mathbb{E}{[q \sim P(Q), o \sim \pi_{\theta_{old}}(O|q)]} \frac{1}{|o|} \sum_{t=1}^{|o|} \min \left[ \frac{\pi_\theta(o_{t} | q, o_{<t})}{\pi_{\theta_{old}}(o_{t} | q, o_{<t})} A_{t}, \text{clip} \left( \frac{\pi_\theta(o_{t} | q, o_{<t})}{\pi_{\theta_{old}}(o_{t} | q, o_{<t})}, 1 - \varepsilon, 1 + \varepsilon \right)  A_{t} \right] ,
\end{equation}

The above equation defines the surrogate loss used in PPO. The expectation is computed with respect to $q$ drawn from $P(Q)$ and $o$ sampled according to the previous policy $\pi_{\theta_{old}}(O|q)$. The loss is averaged across the time steps $t$ for each output sequence of length $|o|$. For each time step $t$, the objective selects the minimum value between the scaled advantage $A_t$ multiplied by the ratio of current to previous policy probabilities, and a version where this ratio is clipped within $[1 - \varepsilon, 1 + \varepsilon]$, thus restricting large policy updates and encouraging stable training dynamics.

Here, $\pi_{\theta}$ denotes the current policy model while $\pi_{\theta_{old}}$ refers to the previous version. The variables $q$ and $o$ correspond to questions and outputs, which are drawn from the question dataset and the outputs generated by the former policy $\pi_{\theta_{old}}$, respectively. The hyper-parameter $\varepsilon$, connected to clipping in PPO, serves to promote stable training. The term $A_t$ signifies the advantage, which is calculated using Generalized Advantage Estimation (GAE) \citep{gae}, leveraging the reward signals $\{r_{\ge t}\}$ and the estimated value function $V_{\psi}$. Therefore, PPO necessitates concurrent training of a value function together with the policy network. To address excessive reward model optimization, a common practice is to introduce a per-token KL divergence penalty in the reward, relative to a reference model, for each token as described in \citep{ouyang2022training}, that is,

\begin{equation}
    r_{t} = r_\varphi(q, o_{\le t}) - \beta \log\frac{\pi_{\theta}(o_{t}|q, o_{<t})}{\pi_{ref}(o_{t}|q, o_{<t})},
\label{eq:PPO-reward}
\end{equation}

Here, $r_\varphi$ represents the reward model, while $\pi_{ref}$ denotes the reference model—commonly the original SFT model. The parameter $\beta$ serves as the weighting factor for the KL divergence penalty.

As the value function employed in PPO is typically another model of comparable size as the policy model, it brings a substantial memory and computational burden. Additionally, during RL training, the value function is treated as a baseline in the calculation of the advantage for variance reduction. While in the LLM context, usually only the last token is assigned a reward score by the reward model, which may complicate the training of a value function that is accurate at each token. To address this, as shown in Figure \ref{fig:grpo}, we propose Group Relative Policy Optimization (GRPO), which obviates the need for additional value function approximation as in PPO, and instead uses the average reward of multiple sampled outputs, produced in response to the same question, as the baseline. More specifically, for each question $q$, GRPO samples a group of outputs $\{o_1, o_2, \cdots, o_G\}$  from the old policy  $\pi_{\theta_{old}}$  and then optimizes the policy model by maximizing the following objective:

\begin{equation}
\footnotesize
\begin{split}
    \mathcal{J}_{GRPO}(\theta) &= \mathbb{E}_{q \sim P(Q),\, \{o_i\}_{i=1}^G \sim \pi_{\theta_{old}}(O|q)}  \\
    & \frac{1}{G} \sum_{i=1}^G \frac{1}{|o_i|} \sum_{t=1}^{|o_i|} \bigg\{ \min \Bigg[ \frac{\pi_\theta(o_{i,t} | q, o_{i,<t})}{\pi_{\theta_{old}}(o_{i,t} | q, o_{i,<t})} \hat{A}_{i,t}, \\
    & \hspace{5em} \text{clip}\Bigg( \frac{\pi_\theta(o_{i,t} | q, o_{i,<t})}{\pi_{\theta_{old}}(o_{i,t} | q, o_{i,<t})}, 1 - \varepsilon, 1 + \varepsilon \Bigg) \hat{A}_{i,t} \Bigg] \\
    & \hspace{2em} - \beta \mathbb{D}_{KL}[\pi_{\theta} \parallel \pi_{ref}] \bigg\} ,
\end{split}
\label{eq:GRPO-obj}
\end{equation}

where $\varepsilon$ and $\beta$ are hyper-parameters, and $\hat{A}_{i,t}$  is the advantage calculated based on relative rewards of the outputs inside each group only, which will be detailed in the following subsections. The group relative way that GRPO leverages to calculate the advantages, aligns well with the comparative nature of rewards models,  as reward models are typically trained on datasets of comparisons between outputs on the same question. Also note that, instead of adding KL penalty in the reward, GRPO regularizes by directly adding the KL divergence between the trained policy and the reference policy to the loss, avoiding complicating the calculation of  $\hat{A}_{i,t}$. And different from the KL penalty term used in (\ref{eq:PPO-reward}), we estimate the KL divergence with the following unbiased estimator \citep{kl_approx}:

\begin{equation}
\small
    \mathbb{D}_{KL}\left[\pi_{\theta} || \pi_{ref}\right] = \frac{\pi_{ref}(o_{i,t}|q,o_{i,<t})}{\pi_{\theta}(o_{i,t}|q,o_{i,<t})} - \log\left(\frac{\pi_{ref}(o_{i,t}|q,o_{i,<t})}{\pi_{\theta}(o_{i,t}|q,o_{i,<t})}\right) - 1,
\end{equation}

which is guaranteed to be positive.

\begin{algorithm}[t]
  \small
  \caption{Iterative Group Relative Policy Optimization}
  \textbf{Input} initial policy model $\pi_{\theta_{\text{init}}}$; reward models $r_\varphi$; task prompts $\mathcal{D}$; 
  hyperparameters $\varepsilon$, $\beta$, $\mu$
  \begin{algorithmic}[1]
    \State Initialize policy model $\pi_\theta$ with $\pi_{\theta_{\text{init}}}$
    \For{iteration = 1, \dots, I}
       \State Set reference model $\pi_{ref}$ to current $\pi_\theta$
      \For{step = 1, \dots, M}
      \State Draw minibatch $\mathcal{D}_b$ from $\mathcal{D}$
      \State Assign $\pi_{\theta_{old}} \leftarrow \pi_{\theta}$
      \State For every prompt $q$ in $\mathcal{D}_b$, generate a set of $G$ candidate outputs $\{o_i\}_{i=1}^G$ by sampling from $\pi_{\theta_{old}} (\cdot \mid q)$
      \State Evaluate corresponding rewards $\{r_i\}_{i=1}^{G}$ for outputs $o_i$ using $r_{\varphi}$
      \State Estimate group relative advantages $\hat{A}_{i,t}$ at each token position $t$ within $o_i$
      \For{GRPO iteration = 1, \dots, $\mu$}
        \State Update $\pi_{\theta}$ by optimizing the GRPO objective defined in Equation \ref{eq:GC-GRPO}
      \EndFor
    \EndFor 
    \State Refine $r_\varphi$ via ongoing training with a replay buffer.
    \EndFor 
  \end{algorithmic}
  \textbf{Output} $\pi_\theta$
  \label{alg:iter-grpo}
\end{algorithm}

\subsubsection{Outcome Supervision RL with GRPO} Specifically, for any given question $q$, a set of $G$ outputs $\{o_1, o_2, \cdots, o_G\}$ is generated using the previous iteration of the policy model, $\pi_{\theta_{old}}$. Each output is evaluated by a reward model, resulting in a set of scores $\mathbf{r}=\{r_1, r_2, \cdots, r_G\}$. These reward values are then standardized by removing the mean and dividing by the standard deviation within the group. Under outcome supervision, the normalized reward associated with each output $o_i$ is assigned as the advantage $\hat{A}_{i, t}$ to all tokens within that output, such that $\hat{A}_{i, t} = \widetilde{r}_i = \frac{r_i- {\rm mean}(\mathbf{r})}{{\rm std}(\mathbf{r})}$. The policy is then updated by maximizing the objective shown in equation (\ref{eq:GRPO-obj}).

\subsubsection{Process Supervision RL with GRPO} Outcome supervision delivers feedback solely after the generation of the entire output, which can be inadequate or inefficient for guiding policies in intricate mathematical scenarios. In line with \cite{wang2023math}, we investigate process supervision, wherein a reward is assigned at the conclusion of each individual reasoning stage. More precisely, given a prompt $q$ and $G$ sampled solutions $\{o_1, o_2, \cdots, o_G\}$, a process reward model assesses and assigns scores to every reasoning step, resulting in a collection of stepwise rewards: $\mathbf{R} = \{ \{r_1^{index(1)},\cdots,r_1^{index(K_1)}\}, \cdots,  \{r_G^{index(1)},\cdots,r_G^{index(K_G)}\} \}$. Here, $index(j)$ denotes the position of the terminal token for the $j$-th reasoning step, and $K_i$ signifies the number of steps present in the $i$-th sample. These raw rewards are further standardized by subtracting the mean and dividing by the standard deviation, i.e., $\widetilde{r}_i^{index(j)} = \frac{r_i^{index(j)} - {\rm mean(\mathbf{R})}}{{\rm std(\mathbf{R})}}$.

The advantage for every token is then computed as the cumulative sum of the normalized rewards across subsequent reasoning steps, expressed as $\hat{A}_{i, t} = \sum_{index(j) \ge t} \widetilde{r}_i^{index(j)}$. Finally, the policy is updated to maximize the objective presented in equation (\ref{eq:GRPO-obj}).

\subsubsection{Iterative RL with GRPO} During reinforcement learning, as the training evolves, earlier versions of the reward model may no longer provide effective guidance for the updated policy model. To address this limitation, we investigate iterative RL using GRPO. As detailed in Algorithm \ref{alg:iter-grpo}, our iterative GRPO framework constructs new datasets for reward model training by sampling trajectories from the current policy model. The reward model undergoes continued training with a replay mechanism, ensuring that 10\% of the training data comes from prior iterations. Subsequently, we update the reference model to match the latest policy model and proceed to refine the policy model using the updated reward model.

\subsection{Training and Evaluating DeepSeekMath-RL}

Our reinforcement learning experiments are carried out using DeepSeekMath-Instruct 7B. The RL training utilizes chain-of-thought-formatted questions from the SFT data that are connected to GSM8K and MATH, totaling approximately 144,000 questions. We deliberately omit additional SFT questions in order to analyze how RL influences benchmarks without training data during the RL period. The reward model training set is built according to \citep{wang2023math}. The reward model is initially trained starting from DeepSeekMath-Base 7B, employing a learning rate of $2 \times 10^{-5}$. Within the GRPO framework, the policy model uses a learning rate of $1 \times 10^{-6}$, and the KL divergence coefficient is fixed at $0.04$. Each individual question yields $64$ sampled completions. Both the maximum sequence length and the training batch size are configured at $1024$. After each phase of exploration, the policy model is updated once. To assess performance, DeepSeekMath-RL 7B is benchmarked in the same manner as DeepSeekMath-Instruct 7B. Here, tasks such as GSM8K and MATH with chain-of-thought reasoning are considered as in-domain, while all other benchmarks are treated as out-of-domain tasks for DeepSeekMath-RL 7B.

Table \ref{tab:sft_rl_math} presents a comparative analysis of open- and closed-source models employing both chain-of-thought and tool-augmented reasoning approaches on English and Chinese evaluation suites. The key observations are as follows: 1) DeepSeekMath-RL 7B achieves 88.2\% accuracy on GSM8K and 51.7\% on MATH using chain-of-thought methods. This sets a new benchmark among open-source models within the 7B to 70B scale and exceeds the performance of most closed-source alternatives. 2) Notably, DeepSeekMath-RL 7B is exclusively fine-tuned on chain-of-thought-formatted instructional data from GSM8K and MATH, beginning with the DeepSeekMath-Instruct 7B checkpoint. Despite its limited training corpus, DeepSeekMath-RL 7B delivers superior results compared to DeepSeekMath-Instruct 7B on all evaluation fronts, highlighting the substantial benefits of reinforcement learning.

\section{Discussion}
This section presents the outcomes derived from both pre-training and reinforcement learning (RL) experiments.

\subsection{Lessons Learnt in Pre-Training}

We begin by detailing our pre-training experiences. Except where indicated, the training configurations follow those described in Section \ref{sec:quality-policy}. It should be highlighted that the term DeepSeekMath Corpus in this context denotes the 89B-token version obtained from the second round of data gathering.

\subsubsection{Code Training Benefits Mathematical Reasoning}

An often-cited but largely untested assumption posits that training on code enhances reasoning capabilities. In this work, we provide initial evidence addressing this claim in the context of mathematics: training with code leads to better mathematical reasoning in models, regardless of whether they are utilizing external tools.

To study how code training affects mathematical reasoning, we experimented with the following two-stage training and one-stage training settings:

\textbf{Two-Stage Training}

\begin{itemize}[topsep=0pt]
    \item \textbf{Code Training for 400B Tokens $\rightarrow$ Math Training for 150B Tokens}: Our approach involves pretraining DeepSeek-LLM 1.3B on 400B tokens of code prior to conducting a secondary phase of training on 150B math tokens.
    \item \textbf{General Training for 400B Tokens $\rightarrow$ Math Training for 150B Tokens}: For comparative purposes, we replace the code tokens with general tokens—drawn from DeepSeek-AI’s extensive general text corpus—in the initial training stage. This variation allows us to examine whether using code tokens specifically confers greater benefits for enhancing mathematical reasoning capabilities than using general textual data.
\end{itemize}

\textbf{One-Stage Training}

\begin{itemize}[topsep=0pt]
    \item \textbf{Math-Only Training on 150B Tokens}: DeepSeek-LLM 1.3B is trained using a dataset containing 150B math-specific tokens;
    \item \textbf{Joint Training with 400B Code Tokens and 150B Math Tokens}: Sequential math training after code-focused training results in diminished coding abilities. We explore if integrating code and math tokens into a unified training phase could enhance math reasoning capabilities while mitigating catastrophic forgetting in code-related knowledge.
\end{itemize}

\paragraph{Results}
The downstream outcomes across various training configurations are presented in Table \ref{tab:code-math} and Table \ref{tab:code-math-general-eval}.

Code training enhances program-aided mathematical reasoning capabilities in both two-stage and one-stage training paradigms. According to the results in Table \ref{tab:code-math}, during the two-stage training process, solely engaging in code training leads to substantial improvements in solving GSM8K and MATH tasks with Python. Introducing math training as a follow-up in the second stage results in additional performance gains. Notably, in the one-stage training scenario, concurrently incorporating both code and math tokens addresses the catastrophic forgetting problem typically observed in two-stage training, while also providing complementary benefits to both coding tasks (see Table \ref{tab:code-math-general-eval}) and program-aided mathematical reasoning tasks (see Table \ref{tab:code-math}).

Training on code also enhances mathematical reasoning abilities even in the absence of tool use. In the two-stage training framework, commencing with code training alone brings about noticeable improvements. This process further increases the effectiveness of later mathematics-specific training, ultimately culminating in superior performance. In contrast, merging code tokens and mathematical tokens into a unified one-stage training approach appears to weaken mathematical reasoning without tool usage. One hypothesis is that DeepSeek-LLM 1.3B, constrained by its smaller size, may not possess sufficient capacity to thoroughly integrate both coding and mathematical information at the same time.

\subsubsection{ArXiv Papers Seem Ineffective in Improving Mathematical Reasoning}

ArXiv papers are commonly included as a component of math pre-training data \citep{minerva,gpt-f,llemma,mathpile}. However, detailed analysis regarding their impact on mathematical reasoning has not been extensively conducted. Perhaps counter-intuitively, according to our experiments, arXiv papers seem ineffective in improving mathematical reasoning. We experiment with models of different sizes, including DeepSeek-LLM 1.3B and DeepSeek-Coder-Base-v1.5 7B \citep{deepseek-coder}, using arXiv corpora that underwent varied processing pipelines:

\begin{itemize}[topsep=0pt]
    \item \textbf{MathPile} \citep{mathpile}: comprises 8.9B tokens, curated using heuristic techniques for data cleaning and filtering, and contains a majority of scientific documents from arXiv (over 85\%);
    \item \textbf{ArXiv-RedPajama} \citep{redpajama}: consists of all arXiv LaTeX sources, excluding preambles, comments, macros, and reference sections, resulting in a dataset of 28.0B tokens.
\end{itemize}

For our experiments, DeepSeek-LLM 1.3B was trained independently on 150B tokens and DeepSeek-Coder-Base-v1.5 7B was trained on 40B tokens using individual arXiv corpora. Results indicate that exposure to arXiv papers offers little benefit for enhancing mathematical reasoning. Training solely on arXiv-based data led both models to show negligible gains or, in some cases, a decline in performance on a range of mathematical evaluation tasks encompassing varying levels of complexity. The evaluation involved quantitative reasoning datasets such as GSM8K and MATH (Table \ref{tab:arxiv-cot}), multiple-choice assessments like MMLU-STEM (Table \ref{tab:arxiv-cot}), and formal mathematical tasks exemplified by miniF2F (Table \ref{tab:arxiv_atp}).

However, this conclusion has its limitations and should be taken with a grain of salt. We have not yet studied:

\begin{itemize}[topsep=0pt]
    \item How arXiv tokens influence mathematics-oriented tasks beyond those explored in this work, for example, transforming formal theorems or proofs into informal language;
    \item The role played by arXiv tokens when integrated with additional data modalities;
    \item If the advantages derived from arXiv papers would persist or change when applied to models of increased size.
\end{itemize}

Thus, further exploration is required, which we leave for future studies.

\subsection{Insights of Reinforcement Learning}
\subsubsection{Towards to a Unified Paradigm} This section establishes a comprehensive framework for examining various training techniques, including SFT, RFT, DPO, PPO, and GRPO. Additionally, we carry out experimental studies to investigate the components influencing this unified approach. Typically, the gradient with respect to the parameter $\theta$ for a given training method is expressed as:

\begin{equation}
    \nabla_{\theta}\mathcal{J}_{\textcolor{red}{\mathcal{A}}}(\theta) = \mathbb{E}[\underbrace{(q,o) \sim \textcolor{red}{\mathcal{D}}}_{Data \ Source}]\left( \frac{1}{|o|} \sum_{t=1}^{|o|}  \underbrace{GC_{{\mathcal{A}}}(q, o, t, \textcolor{red}{\pi_{{rf}}})}_{Gradient \ Coefficient}  \nabla_{\theta}\log \pi_{\theta}(o_t | q, o_{<t})\right).
\label{eq:objective}
\end{equation}

There exist three key components: 1) \textit{Data Source $\mathcal{D}$}, which determines the training data; 2) \textit{Reward Function $\pi_{{rf}}$}, which is the source of the training reward signal; 3) \textit{Algorithm $\mathcal{A}$}: which processes the training data and the reward signal to the gradient coefficient $GC$ that determines the magnitude of the penalty or reinforcement for the data. We analyze several representative methods based on such a unified paradigm:

\begin{itemize}
    \item  \textbf{Supervised Fine-tuning (SFT)}: In SFT, a pretrained model undergoes additional training on a dataset curated by humans for supervised fine-tuning.
    \item \textbf{Rejection Sampling Fine-tuning (RFT)}: RFT continues the fine-tuning process for an SFT model using a subset of outputs generated in response to SFT questions. These outputs are filtered based on the correctness of their answers before being used for further training.
    \item \textbf{Direct Preference Optimization (DPO)}: DPO enhances the SFT model by fine-tuning it with additional outputs sampled from the SFT model itself, applying a pair-wise DPO loss to optimize preferences directly.
    \item \textbf{Online Rejection Sampling Fine-tuning (Online RFT)}: Unlike standard RFT, Online RFT initializes the policy model from the SFT model but updates the model in an online manner, utilizing augmented outputs drawn in real time from the evolving policy model for further training.
    \item \textbf{PPO/GRPO}: With PPO/GRPO, the policy model is initialized using the SFT model and then further trained using samples generated by the current policy model to reinforce its behavior.
\end{itemize}

The elements constituting these approaches are outlined in Table \ref{tab:data_policy}. For an in-depth explanation of the derivation procedure, see Appendix \ref{app:analysis-rl}.

\paragraph{Observation about Data Source} We categorize data sources into two types: online sampling and offline sampling. Online sampling refers to collecting training data generated by the ongoing, active training policy model. In contrast, offline sampling pertains to data drawn from the initial SFT model's outputs. Methods such as RFT and DPO utilize an offline approach, whereas Online RFT and GRPO adopt an online approach.

As illustrated in Figure \ref{fig:rl-analysis}, our results indicate that Online RFT achieves markedly better performance than RFT on both benchmarks. In particular, Online RFT matches RFT’s performance during the early phases of training but subsequently attains a clear lead as training progresses, highlighting the benefits of the online training regime. This outcome is to be expected, since at the outset, the actor and SFT model are highly similar, resulting in the collected samples reflecting only minor distinctions. As training advances, the actor's sampled data diverges more significantly, making the real-time sampling process increasingly beneficial.

\paragraph{Observation about Gradient Coefficient} In our methodology, the reward signal informs the gradient coefficient, which in turn directs the model parameter updates. Within our experimental setup, the reward mechanism is delineated into two categories: `Rule' and `Model'. The `Rule' category entails assessment of response quality by verifying the accuracy of answers, while the `Model' category involves the use of a trained reward model that assigns a score to each response, with this model being trained on data annotated using rule-based evaluations. A salient distinction between GRPO and Online RFT emerges from Equations \ref{eq:GC-RFT} and \ref{eq:GC-GRPO}: GRPO modifies its gradient coefficient in accordance with the reward magnitude assigned by the reward model. This adaptive process enables varying degrees of encouragement or discouragement depending on how favorable the response is. Conversely, Online RFT omits this adaptive aspect, as it neither penalizes incorrect outputs nor differentiates among correct answers, treating all correct responses identically in terms of reinforcement strength.

As shown in Figure \ref{fig:rl-analysis}, GRPO outperforms online RFT, underscoring the effectiveness of adjusting positive and negative gradient coefficients. Moreover, the results demonstrate that GRPO+PS achieves better outcomes than GRPO+OS, reflecting the advantages of employing step-aware, fine-grained gradient coefficients. We also investigate the impact of iterative RL by conducting two sequential iterations in our experiments. According to Figure \ref{fig:iter-rl}, the use of iterative RL leads to substantial performance gains, with the most pronounced improvement observed during the initial iteration.

\subsubsection{Why RL Works?} In this study, we utilize reinforcement learning on a selected subset of instruction tuning data, resulting in notable improvements over the baseline instruction tuning model. To shed light on the mechanisms underlying these gains, we assess Pass@K and Maj@K accuracies for both the Instruct and RL-trained models across two different benchmarks. As illustrated in Figure \ref{fig:maj-pass}, while RL drives improvements in Maj@K performance, it does not significantly affect Pass@K. This suggests that reinforcement learning primarily enhances the resilience of the output distribution—specifically, \textbf{the observed gains appear to stem from an increased likelihood of surfacing the correct answer among the TopK predictions, rather than from an intrinsic advancement in core abilities.} In line with our findings, \citep{wang2023making} discuss a \textbf{misalignment problem} in reasoning tasks encountered with SFT models, demonstrating that reasoning efficacy can be enhanced via a range of preference alignment approaches \citep{yuan2023rrhf,song2023preference,wang2023making}.

\subsubsection{How to Achieve More Effective RL?}
\label{sec:effec_rl}
Our results show that RL exhibits strong performance on mathematical reasoning challenges. Moreover, we propose a unified framework that brings together various influential training strategies. In this framework, each approach can be interpreted as employing RL either in its original or in a more abstracted form. As outlined in Equation \ref{eq:objective}, the essential elements include: Data Source, Algorithm, and Reward Function. We outline several possible avenues for advancement concerning these three aspects.

\paragraph{Data Source} The foundational input for all training methodologies is the data source. Within the RL framework, we define this data source as unlabeled questions along with their corresponding outputs generated by the policy model. In this study, we restrict ourselves to questions utilized during the instruction tuning phase and employ a basic nucleus sampling approach to produce responses. This constraint might explain why our RL framework yields improvements solely in the Maj@K metric. Moving forward, we plan to apply our RL framework to question prompts outside the training distribution, pairing this with \textbf{advanced sampling (decoding) strategies} such as those utilizing tree-search techniques \citep{yao2023tree}. Moreover, \textbf{efficient inference techniques} \citep{xia-etal-2023-speculative,leviathan2023fast,kwon2023efficient,xia2024unlocking}, which significantly impact the efficiency with which policy models explore, are also expected to have a critical influence.

\paragraph{Algorithms} Algorithms utilize both data and the reward signal to determine the gradient coefficient used for updating model parameters. Referring to Equation \ref{eq:objective}, most existing methods effectively place complete \textbf{TRUST} in the reward function's signal when adjusting the conditional probability associated with specific tokens. Nonetheless, in highly intricate tasks, there is no guarantee regarding the dependability of the reward signal. As a case in point, even the extensively curated PRM800K datasets \citep{lightman2023let}, which were annotated by skilled annotators, still exhibit an error rate of nearly 20\% in their labels\footnote{\url{https://github.com/openai/prm800k/issues/12\#issuecomment-1728491852}}. Given this, our aim is to investigate reinforcement learning algorithms that can maintain performance in the presence of noisy reward signals. We posit that alignment strategies such as \textbf{WEAK-TO-STRONG} \citep{burns2023weak} have the potential to significantly transform current learning paradigms.

\paragraph{Reward Function} The reward function serves as the fundamental provider of the training signal. Within the context of RL, this function is typically represented by a neural reward model. We identify three pivotal areas of exploration for improving reward models: 1) \textbf{Improving the generalization of the reward model.} It is crucial that the reward model demonstrates strong generalization, enabling it to respond robustly to out-of-distribution queries and sophisticated decoded outputs. Absent such generalizability, reinforcement learning may simply maintain the existing distribution of LLMs rather than fostering core enhancements; 2) \textbf{Capturing the uncertainty in the reward model.} Accounting for the reward model’s uncertainty can facilitate interaction between suboptimal reward models and algorithms that transfer from weak to strong learning; 3) \textbf{Developing efficient and high-quality process reward models} capable of yielding detailed, step-level training feedback throughout the reasoning pipeline \citep{lightman2023let,wang2023math}.

\section{Conclusion, Limitation, and Future Work} 

In this work, we introduce DeepSeekMath, a model that surpasses existing open-source counterparts on the MATH benchmark at competition level and nearly matches the results achieved by proprietary systems. The model is built upon DeepSeek-Coder-v1.5 7B, followed by continual pretraining over 500B tokens, including a substantial subset of 120B mathematics-specific tokens extracted from Common Crawl. Our comprehensive ablation analyses reveal that web sources contain considerable high-quality mathematical data, whereas contributions from arXiv were less advantageous than anticipated. We propose Group Relative Policy Optimization (GRPO), an adaptation of Proximal Policy Optimization (PPO), which enhances mathematical reasoning performance while reducing memory requirements. Experimental findings indicate that GRPO delivers substantial benefits, even when DeepSeekMath-Instruct 7B already achieves strong benchmark scores. Additionally, we outline a cohesive framework for understanding a spectrum of current techniques and discuss promising avenues for advancing reinforcement learning methodologies.

While DeepSeekMath~delivers strong results on benchmarks that emphasize quantitative reasoning, its performance in geometry and theorem-proving tasks lags behind that of proprietary models. For example, our preliminary tests revealed that the model struggles with geometry questions involving triangles and ellipses, suggesting a possible bias in the datasets used during pre-training and fine-tuning. Furthermore, owing to the limitations imposed by its scale, DeepSeekMath~does not match GPT-4 in terms of few-shot learning capabilities. Unlike GPT-4, which benefits from few-shot inputs by exhibiting improved results, DeepSeekMath~performs comparably in both zero-shot and few-shot scenarios. Moving forward, we plan to refine our engineered data selection process to curate a higher-quality pre-training corpus. Additionally, we intend to investigate new strategies for more efficient reinforcement learning of LLMs, as discussed in Section \ref{sec:effec_rl}.

\end{document}